\usetikzlibrary{arrows.meta,positioning}
\definecolor{probColor}{RGB}{180,60,60}
\definecolor{solColor}{RGB}{70,110,190}
\definecolor{coreColor}{RGB}{40,130,90}

\begin{tikzpicture}[
    % 定义通用节点样式
    node distance=1.5cm and 0.8cm, % 节点间的距离 (纵向 and 横向)
    % 1. 顶层矩形样式
    topBox/.style={
        rectangle, 
        draw=probColor!80, 
        fill=probColor!10, 
        very thick,
        rounded corners,
        minimum width=12cm, % 设定宽度以覆盖下方的立柱
        minimum height=1.2cm,
        align=center,
        font=\bfseries\large,
        inner sep=5pt
    },
    % 2. 中间立柱样式
    pillarBox/.style={
        rectangle, 
        draw=solColor!80, 
        fill=solColor!5, 
        thick,
        minimum width=3.8cm, % 三个立柱平分宽度
        minimum height=4cm,  % 立柱高度
        align=center,
        text width=3.5cm,    % 限制文本宽度以自动换行
        inner sep=5pt,
        font=\small
    },
    % 3. 底层横向矩形样式
    bottomBox/.style={
        rectangle, 
        draw=coreColor!80, 
        fill=coreColor!10, 
        thick,
        rounded corners,
        minimum width=8cm,
        minimum height=1cm,
        align=center,
        font=\bfseries
    },
    % 4. 箭头样式
    arrow/.style={
        -{Stealth[scale=1.2]}, % 使用专业的 Stealth 箭头
        thick,
        draw=gray!80
    },
    % 特殊样式：带有“解耦”文字的加粗箭头
    decoupleArrow/.style={
        -{Stealth[scale=1.5]},
        ultra thick,
        draw=coreColor,
        shorten >=2pt,
        shorten <=2pt
    }
]

% ==========================================================
% 第一步：放置节点
% ==========================================================

% 1. 顶层：总的问题
\node (topnode) [topBox] {病态系统 (Ill-conditioned System)};

% 2. 中间层：三个分立的研究问题 (立柱)
% 先放中间的，再放左右的
\node (prob2) [pillarBox, below=of topnode] {\textbf{问题二}\\[2ex]包含恶劣单元的\\有限元网格计算};
\node (prob1) [pillarBox, left=of prob2] {\textbf{问题一}\\[2ex]不可拉伸弹性绳\\的动力学仿真};
\node (prob3) [pillarBox, right=of prob2] {\textbf{问题三}\\[2ex]使用CB-CMS局部特征基函数来求解的\\广义特征问题};

% 3. 底层：核心思路 (总)
% 计算三个立柱的中心对齐位置
\node (coupling) [bottomBox, below=1.2cm of prob2] {存在软硬耦合的基函数};
\node (transform) [bottomBox, below=1.8cm of coupling] {基函数变换 (Basis Function Transformation)};

% ==========================================================
% 第二步：绘制连接箭头
% ==========================================================

% 1. 顶层连向三个立柱
\draw [arrow] (topnode.south) -- (prob1.north);
\draw [arrow] (topnode.south) -- (prob2.north);
\draw [arrow] (topnode.south) -- (prob3.north);

% 2. 三个立柱连向底层思路
% 为了美观，使用汇聚形式，让 prob1 和 prob3 弯曲连向中间位置
\draw [arrow] (prob1.south) -- ++(0,-0.6cm) -| (coupling.north);
\draw [arrow] (prob2.south) -- (coupling.north);
\draw [arrow] (prob3.south) -- ++(0,-0.6cm) -| (coupling.north);

% 3. 思路一连向思路二，旁边写着“解耦”
\draw [decoupleArrow] (coupling.south) -- (transform.north)
    node[midway, right=0.2cm, text=coreColor, font=\bfseries\large] {解耦};

\end{tikzpicture}